\exercisegroup

\begin{enumerate}
\def\labelenumi{\arabic{enumi}.}
\tightlist
\item
  Let S be the set of signs P, P-, X, X-, \(x_1, \ldots, x_n\) , the letters \(x_i\) being of weight 0, P and P- of weight 1, and X, X- of weight 2. For each word A in L(S) we define the variance of A as follows: each letter \(x_i\) , and the signs P, X, have variance 0; the signs P-, X- have variance 1; the variance of A is the sum * (in the field F₂) * of the variances of the signs which appear in A (in other words, it is 0 if A contains an even number of signs of variance 1, and is 1 otherwise).
\end{enumerate}

A signed echelon type is a balanced word A in L(S) (Chapter I, Appendix, no. 3) which satisfies one of the following two conditions: (1) A is one of the letters \(x_i\) ; (2) if A is of the form \(f\mathbf{A}_1\mathbf{A}_2 \ldots \mathbf{A}_p\) (Chapter I, Appendix, no. 3, Corollary 2 of Proposition 2), where \(f\) is one of the signs \(\mathsf{P}, \mathsf{P}^{-}, \mathsf{X}, \mathsf{X}^{-}, p\) is equal to 1 or 2 and the \(\mathbf{A}_i\) ( \(1 \leqslant i \leqslant p\) ) are balanced words, then each of the \(\mathbf{A}_i\) is a signed echelon type; further, if \(f = \mathsf{X}\) , then \(\mathbf{A}_1\) and \(\mathbf{A}_2\) have variance 0, and if \(f = \mathsf{X}^{-}\) , then \(\mathbf{A}_1\) and \(\mathbf{A}_2\) have variance 1.

A signed echelon type is said to be covariant if it has variance 0, contravariant if it has variance 1. If in a signed echelon type A we replace P- and X- by P and X, respectively, we obtain an echelon type A* (§ 1, Exercise 1). Every realization of the echelon type A* on n terms E₁, \ldots, Eₙ is called a realization of the signed echelon type A on E₁, \ldots, Eₙ, and is written A(E₁, \ldots, Eₙ).

Let \(\mathbf{E}_1, \ldots, \mathbf{E}_n, \mathbf{E}_1', \ldots, \mathbf{E}_n'\) be sets, and let \(f_i\) be a mapping of \(\mathbf{E}_i\) into \(\mathbf{E}_i'\) ( \(1 \leqslant i \leqslant n\) ). Show that to each signed echelon type S on \(x_1, \ldots, x_n\) we may associate a mapping \(\{f_1, \ldots, f_n\}^{S}\) having the following properties:

\begin{enumerate}
\def\labelenumi{(\roman{enumi})}
\item
  If S is covariant (resp. contravariant), then \(\{f_1, \ldots, f_n\}^{S}\) is a mapping of \(\mathrm{S}(\mathrm{E}_1, \ldots, \mathrm{E}_n)\) into \(\mathrm{S}(\mathrm{E}_1', \ldots, \mathrm{E}_n')\) (resp. of \(\mathrm{S}(\mathrm{E}_1', \ldots, \mathrm{E}_n')\) into \(\mathrm{S}(\mathrm{E}_1, \ldots, \mathrm{E}_n)\) ).
\item
  If S is a letter \(x_{i}\) , then \(\{f_1, \ldots, f_n\}^{S}\) is \(f_{i}\) .
\item
  If S is PT (resp. P-T), and if \(g = \{f_1, \ldots, f_n\}^{\mathrm{T}}\) is a mapping of F into \(\mathbf{F}'\) , then \(\{f_1, \ldots, f_n\}^{S}\) is the extension of \(g\) to sets of subsets (resp. the extension of \(\overline{g}^1\) to sets of subsets).
\item
  If S is XTU or X-TU, if \(\{f_{1},\ldots,f_{n}\}^{T}\) is a mapping \(g:F\to F'\) and if \(\{f_{1},\ldots,f_{n}\}^{U}\) is a mapping \(h:G\to G'\) , then \(\{f_{1},\ldots,f_{n}\}^{S}\) is the extension \(g\times h:F\times G\to F'\times G'\) .
\end{enumerate}

The mapping \(\{f_1, \ldots, f_n\}^{S}\) is called the signed canonical extension of \(f_1, \ldots, f_n\) corresponding to the signed echelon type S. If S is an echelon type (i.e., if \(\mathbb{P}^{-}\) and \(\mathbb{X}^{-}\) do not appear in S), then the signed canonical extension \(\{f_1, \ldots, f_n\}^{S}\) is equal to \(\langle f_1, \ldots, f_n \rangle^{S}\) .

Show that, if \(f_{i}\) is a mapping of \(\mathbf{E}_i\) into \(\mathbf{E}_i'\) , and if \(f_{i}'\) is a mapping of \(\mathbf{E}_i'\) into \(\mathbf{E}_i''\) ( \(1 \leqslant i \leqslant n\) ), then for a covariant signed echelon type S we have

\[
\left\{f _ {1} ^ {\prime} \circ f _ {1}, \dots , f _ {n} ^ {\prime} \circ f _ {n} \right\} ^ {\mathrm{S}} = \left\{f _ {1} ^ {\prime}, \dots , f _ {n} ^ {\prime} \right\} ^ {\mathrm{S}} \circ \left\{f _ {1}, \dots , f _ {n} \right\} ^ {\mathrm{S}},
\]

and for a contravariant signed echelon type we have

\[
\left\{f _ {1} ^ {\prime} \circ f _ {1}, \dots , f _ {n} ^ {\prime} \circ f _ {n} \right\} ^ {\mathrm{S}} = \left\{f _ {1}, \dots , f _ {n} \right\} ^ {\mathrm{S}} \circ \left\{f _ {1} ^ {\prime}, \dots , f _ {n} ^ {\prime} \right\} ^ {\mathrm{S}}.
\]

Deduce, that if \(f_i\) is a bijection of \(\mathbf{E}_i\) onto \(\mathbf{E}_i'\) and if \(f_i'\) is its inverse bijection ( \(1 \leqslant i \leqslant n\) ), then \(\{f_1, \ldots, f_n\}^S\) is a bijection and \(\{f_1', \ldots, f_n'\}^S\) its inverse. Moreover, if \(S^*\) is the (unsigned) echelon type corresponding to the signed echelon type \(S\) , then \(\{f_1, \ldots, f_n\}^S\) is equal to \(\langle f_1, \ldots, f_n \rangle^{S^*}\) or to \(\langle f_1', \ldots, f_n' \rangle^{S^*}\) according as \(S\) is covariant or contravariant.

\begin{enumerate}
\def\labelenumi{\arabic{enumi}.}
\setcounter{enumi}{1}
\tightlist
\item
  Let S be a signed echelon type on \(n + m\) letters (Exercise 1). Let \(\Sigma\) be a species of structures having \(x_{1}, \ldots, x_{n}\) as principal base sets and \(A_{1}, \ldots, A_{m}\) as auxiliary base sets, whose typical characterization is of the form \(s \in \mathfrak{P}(S(x_{1}, \ldots, x_{n}, A_{1}, \ldots, A_{m}))\) . Show that we may define a notion of \(\sigma\) -morphism for this species of structures as follows: given n sets \(E_{1}, \ldots, E_{n}\) endowed with a structure U of species \(\Sigma\), n sets \(E_{1}^{\prime}, \ldots, E_{n}^{\prime}\) endowed with a structure \(U^{\prime}\) of species \(\Sigma\) , and mappings \(f_{i}: E_{i} \to E_{i}^{\prime} (1 \leqslant i \leqslant n)\) , then \((f_{1}, \ldots, f_{n})\) is said to be a \(\sigma\) -morphism if the mappings \(f_{i}\) satisfy the following conditions:
\end{enumerate}

\begin{enumerate}
\def\labelenumi{(\roman{enumi})}
\tightlist
\item
  If S is a covariant echelon type, then
\end{enumerate}

\[
\{f _ {1}, \dots , f _ {n}, I _ {1}, \dots , I _ {m} \} ^ {S} \langle U \rangle \subset U ^ {\prime}.
\]

\begin{enumerate}
\def\labelenumi{(\roman{enumi})}
\setcounter{enumi}{1}
\tightlist
\item
  If \(S\) is a contravariant echelon type, then
\end{enumerate}

\[
\left\{f _ {1}, \dots , f _ {n}, I _ {1}, \dots , I _ {m} \right\} ^ {S} \langle U ^ {\prime} \rangle \subset U.
\]

Show that by assigning the variances suitably we may obtain in this way the definition of the morphisms for order structures, algebraic structures, and topological structures.

\begin{enumerate}
\def\labelenumi{\arabic{enumi}.}
\setcounter{enumi}{2}
\item
  Let A, B, C be three sets endowed with structures of the same species \(\Sigma\) , let \(f\) be a morphism of A onto B, and let \(g\) be a morphism of B into C. Show that if \(g \circ f\) is an isomorphism of A onto C, then \(g\) and \(f\) are isomorphisms.
\item
  Let A, B, C, D be four sets endowed with structures of the same species \(\Sigma\) , let \(f\) be a morphism of A into B, \(g\) a morphism of B into C, and \(h\) a morphism of C into D. Show that if \(g \circ f\) and \(h \circ g\) are isomorphisms, then \(f\) , \(g\) , and \(h\) are isomorphisms (cf.~Chapter II, §3, Exercise 9).
\item
  Let A, B be two sets endowed with structures \(\mathcal{S}\) , \(\mathcal{S}'\) of the same species \(\Sigma\) . Let \(f\) be a morphism of A into B, and let \(g\) be a morphism of B into A. Let M (resp. N) be the set of all \(x \in A\) (resp. \(y \in B\) ) such that \(g(f(x)) = x\) (resp. \(f(g(y)) = y\) ). Suppose that \(\mathcal{S}\) (resp. \(\mathcal{S}'\) ) induces on M (resp. N) a structure of species \(\Sigma\) . Show that M and N, endowed with these structures, are isomorphic.
\item
  Let \(\Sigma\) be the species of structure having a principal base set A and an auxiliary set k, whose generic structure (F, H) has as typical characterization
\end{enumerate}

\[
\mathbf {F} \in \mathfrak {P} ((\mathbf {A} \times \mathbf {A}) \times \mathbf {A}) \times \mathfrak {P} ((\mathbf {A} \times \mathbf {A}) \times \mathbf {A}) \times \mathfrak {P} ((k \times \mathbf {A}) \times \mathbf {A}) \quad \text { and } \quad \mathbf {H} \in \mathfrak {P} (\mathbf {A})
\]

and whose axiom (which can be written more explicitly) is the following: ``F is a structure of commutative \(k\) -algebra with identity element, and H is an irreducible ideal of the algebra \(\mathbf{A}''\) (we recall that an irreducible ideal in a \(k\) -algebra A is an ideal \(\mathfrak{a}\) such that, for any ideals \(\mathfrak{b}, \mathfrak{c}\) of A satisfying \(\mathfrak{b} \cap \mathfrak{c} = \mathfrak{a}\) , either \(\mathfrak{b} = \mathfrak{a}\) or \(\mathfrak{c} = \mathfrak{a}\) ). If A, A' are two sets endowed respectively with structures (F, H), (F', H') of species \(\Sigma\) , we define the \(\sigma\) -morphisms of A into A' to be homomorphisms \(f\) of \(k\) -algebras, which map identity element to identity element, and satisfy \(f(H) \subset H'\) . Give an example of a family ( \(\mathcal{G}_{\lambda}\) ) of structures of species \(\Sigma\) on a set A, such that there exists a least upper bound for this family of structures (in the ordered set of structures of species \(\Sigma\) on A) but such that this least upper bound is not an initial structure for the family (A \(_{\lambda}\) , \(\mathcal{G}_{\lambda}\) , Id \(_{\lambda}\) ), where A \(_{\lambda}\) is the set A endowed with the structure \(\mathcal{G}_{\lambda}\) , and Id \(_{\lambda}\) is the identity mapping A → A \(_{\lambda}\) . (Consider a polynomial ring A = k{[}T{]}: the irreducible ideals in A are powers of maximal ideals. If F is the \(k\) -algebra structure of A, consider the two structures (F, p \(_1\) ) and (F, p \(_2\) ) of species \(\Sigma\) , where p \(_1\) , p \(_2\) are distinct maximal ideals of A; show that the least upper bound of these two structures is (F, (0)), but that this is not the initial structure for the family in question. To prove this, consider the \(k\) -subalgebra B = k + (p \(_1\) ∩ p \(_2\) ) of A, in which p \(_1\) ∩ p \(_2\) is a maximal ideal, and the injection B → A.) Also give an example of a species of structure \(\Sigma'\) , and a family ( \(\mathcal{G}_{\lambda}'\) ) of structures of species \(\Sigma'\) on a set A', such that the greatest lower bound of this family of structures exists but is not the final structure for the family. (On the dual of A (considered as a vector space) consider the structure of linearly compact \(k\) -coalgebra deduced from the \(k\) -algebra structure of A by transposition, and define the species of structures \(\Sigma'\) by transposition.) * 7. Let \(\Sigma\) be the species of structures which has a principal base set A, and the set R as auxiliary base set, whose generic structure consists of two letters (V, φ) with the typical characterization

\[
\mathbf {V} \in \mathfrak {P} (\mathfrak {P} (\mathbf {A})) \quad \text { and } \quad \varphi \in \mathfrak {P} (\mathbf {R} \times \mathbf {A}),
\]

and whose axioms are (i) the axiom R\{V\} of the species of topological structures (§ 1, no. 4, Example 3), (ii) the relation ``there exists \(a>0\) such that φ is the graph of a continuous injective mapping (with respect to the topology V) of the interval {[}0, a{]} into A''.

If A, \(A'\) are two sets endowed respectively with structures (V, \(\varphi\) ), (V', \(\varphi'\) ) of the species \(\Sigma\) , the \(\sigma\) -morphisms of A into \(A'\) are defined to be the mappings \(f\) of \(\mathbf{A}\) into \(\mathbf{A}'\) which are continuous (with respect to the topologies \(\mathbf{V}, \mathbf{V}'\) ) and whose graph \(\mathbf{F}\) is such that \(\mathbf{F} \circ \varphi \subset \varphi'\) . Show that this notion of \(\sigma\) -morphisms can be defined by the procedure of Exercise 2, and that there exists a product structure on the product of two sets \(\mathbf{A}_1, \mathbf{A}_2\) endowed with arbitrary structures of species \(\Sigma\) . But give an example in which the direct image, under the first projection \(\mathbf{pr}_1\) , of the product structure on \(\mathbf{A}_1 \times \mathbf{A}_2\) is not the original structure given on \(\mathbf{A}_1\) (take \(\mathbf{A}_1\) to be a space homeomorphic to \(\mathbf{R}\) ). *

\begin{enumerate}
\def\labelenumi{\arabic{enumi}.}
\setcounter{enumi}{7}
\tightlist
\item
  Let \(\Sigma\) be the species of structures which has a single (principal) base set, whose generic structure \((V, a, b)\) consists of three letters, with the typical characterization
\end{enumerate}

\[
\mathrm{V} \in \mathfrak {P} (\mathfrak {P} (\mathrm{A})) \quad \text { and } \quad a \in \mathrm{A} \quad \text { and } \quad b \in \mathrm{A}
\]

and whose axiom is the relation

\[
\mathrm{R} \{\mathrm{V} \} \quad \text { and } \quad a \neq b,
\]

where \(R\{V\}\) is the axiom of the species of topological structures ( \(\S\) 1, no. 4, Example 3). If A, \(A'\) are two sets endowed respectively with structures (V, a, b), ( \(V'\) , \(a'\) , \(b'\) ) of species \(\Sigma\) , the \(\sigma\) -morphisms of A into \(A'\) are defined to be the continuous mappings \(f: A \to A'\) (with respect to the topologies V, \(V'\) ) such that \(f(a) = a'\) and \(f(b) = b'\) . Show that by replacing \(\Sigma\) by an equivalent species of structures, we can obtain this notion of \(\sigma\) -morphisms by the procedure of Exercise 2. Show that if A, B are two sets endowed with structures S, \(S'\) of species \(\Sigma\) , there exists a product structure on \(A \times B\) , and that the direct image of this structure under \(pr_{1}\) (resp. \(pr_{2}\) ) is S (resp. \(S'\) ). * But give an example where there is no section of \(pr_{1}\) which is a \(\sigma\) -morphism of A into \(A \times B\) (take A to be connected and B discrete). *

* 9. Let \(\Sigma\) be the species of division ring structures. Show that a notion of \(\sigma\) -morphisms for this species of structures may be defined by taking the morphisms of a division ring \(\mathbf{K}\) into a division ring \(\mathbf{K}'\) to be the homomorphisms \(f\) of \(\mathbf{K}\) into \(\mathbf{K}'\) in the sense of Example 2 of no. 4, together with the mapping \(f_0\) of \(\mathbf{K}\) into \(\mathbf{K}'\) for which \(f_0(0) = 0\) and \(f_0(x) = 1\) for all \(x \neq 0\) in \(\mathbf{K}\) . Show that, with this notion of morphisms, we have the following properties: for every division ring \(\mathbf{K}\) (of any characteristic) the division ring structure of \(\mathbf{K}\) induces a division ring structure (isomorphic to that of \(\mathbf{F}_2\) ) on the set \(\{0, 1\}\) ; and that if \(\mathbf{R}\) is the equivalence relation whose equivalence classes are \(\{0\}\) and \(\mathbf{K}^* = \mathbf{K} - \{0\}\) , there exists a quotient structure (isomorphic to that of \(\mathbf{F}_2\) ) of the structure of \(\mathbf{K}\) by the relation \(\mathbf{R}\) .

*10. Let \(\Sigma\) be the species of structures of a complete Archimedean ordered field. For each set A endowed with a structure of species \(\Sigma\) , let \(\varphi_{\mathbf{A}}\) be the unique isomorphism of \(\mathbf{A}\) onto \(\mathbf{R}\) . If \(\mathbf{A}, \mathbf{B}\) are two sets endowed with structures of species \(\Sigma\) , show that the morphisms of \(\mathbf{A}\) into \(\mathbf{B}\) may be taken to be the mappings \(f\) of \(\mathbf{A}\) into \(\mathbf{B}\) such that \(\varphi_{\mathbf{B}}(f(x)) \geqslant \varphi_{\mathbf{A}}(x)\) for all \(x \in \mathbf{A}\) . For this notion of morphisms show that there exist bijective morphisms which are not isomorphisms, although the species \(\Sigma\) is univalent. *

\begin{enumerate}
\def\labelenumi{\arabic{enumi}.}
\setcounter{enumi}{10}
\tightlist
\item
  Let \(\Sigma\) be a species of structures (in a theory \(\mathfrak{G}\) ) which has only one base set. Let \(s \in \mathrm{F}(x)\) be the typical characterization, and \(\mathbf{R}\{x, s\}\) the axiom of \(\Sigma\) . Let \(\mathbf{A}(x)\) denote the set of structures of species \(\Sigma\) on \(x\) . Let \(\sigma\{x, y, s, t\}\) be a term which satisfies the conditions \((\mathrm{MO}_{\mathbf{I}}), (\mathrm{MO}_{\mathbf{II}})\) , and the following condition:
\end{enumerate}

\((MO_{\mathrm{III}}^{\prime})\) Given two sets \(E, E^{\prime}\) (in a theory \(C^{\prime}\) which is stronger than \(C\) ) endowed with structures \(S, S^{\prime}\) of species \(\Sigma\) , respectively, if \(f\) is any isomorphism of \(E\) onto \(E^{\prime}\) , then \(f \in \sigma \{E, E^{\prime}, S, S^{\prime}\}\) .

\begin{enumerate}
\def\labelenumi{(\alph{enumi})}
\tightlist
\item
  Let \(I_{x}\) be the identity mapping of x onto itself. Show that the relation \(Q\{x, s, t\}\) :
\end{enumerate}

\[
s \in \mathrm{A} (x) \quad \text { and } \quad t \in \mathrm{A} (x) \quad \text { and } \quad \mathrm{I} _ {x} \in \sigma \left\{x, x, s, t \right\} \cap \sigma \left\{x, x, t, s \right\}
\]

is an equivalence relation between \(s\) and \(t\) on \(\mathbf{A}(x)\) . Let \(\mathbf{B}(x)\) be the quotient set \(\mathbf{A}(x)/\mathbf{Q}\) , and let \(\varphi_x\) be the canonical mapping of \(\mathbf{A}(x)\) onto \(\mathbf{B}(x)\) . Suppose that the relation \(s' \in \mathbf{B}(x)\) is transportable (see the Appendix for conditions which will ensure that this is so), and let \(\Theta\) denote the species of structures with typical characterization \(s' \in \mathfrak{P}(\mathbf{F}(x))\) and axiom \(s' \in \mathbf{B}(x)\) .

\begin{enumerate}
\def\labelenumi{(\alph{enumi})}
\setcounter{enumi}{1}
\tightlist
\item
  Let \(\overline{\sigma}\{x, y, s', t'\}\) be the set of mappings \(f \in \mathcal{F}(x, y)\) which satisfy the following relation: '' \(s' \in \mathrm{B}(x)\) and \(t' \in \mathrm{B}(y)\) and there exist \(s \in \mathrm{A}(x)\) and \(t \in \mathrm{A}(y)\) such that \(s' = \varphi_x(s)\) and \(t' = \varphi_y(t)\) and \(f \in \sigma\{x, y, s, t\}\) ``. Show that, for the species of structures \(\Theta\) , the term \(\overline{\sigma}\) satisfies \((\mathrm{MO}_{\mathbf{I}})\) , \((\mathrm{MO}_{\mathbf{II}})\) , and \((\mathrm{MO}_{\mathbf{III}})\) and that we have
\end{enumerate}

\[
\sigma \left\{x, y, s, t \right\} \subset \overline {{\sigma}} \left\{x, y, \varphi_ {x} (s), \varphi_ {y} (t) \right\}.
\]

